\documentclass[11pt,a4paper]{article}
\usepackage[T1]{fontenc}
\usepackage{lmodern,amsmath,amssymb,amsthm,booktabs,microtype}
\usepackage[margin=25mm]{geometry}
\usepackage{hyperref}
\hypersetup{colorlinks=true,linkcolor=blue,urlcolor=blue}
\newcommand{\diag}{\operatorname{diag}}
\newcommand{\off}{\operatorname{off}}
\newcommand{\normrows}{\operatorname{normrows}}
\newcommand{\pos}[1]{\left[#1\right]_+}
\newtheorem{proposition}{Proposition}
\setlength{\parindent}{0pt}
\setlength{\parskip}{5pt}
\setlength{\emergencystretch}{3em}
\title{Why these scores belong together in BundleFuse\\
\large Exact relations, information retained, and the final optimization}
\author{}\date{}
\begin{document}
\maketitle

\textbf{The claim developed here.}
The score distinguishes three different questions: whether a candidate occurred
with the observed items, whether it occurred with several of them in the same
historical bundle, and whether it has a similar pattern of relations or attributes
when direct observations are missing. These questions are different mathematical
operations. Their outputs become comparable through catalogue-relative scaling.
The final score is a regularized compromise whose reliance on the combined
evidence increases with the candidate's historical support.

This is an explanation of the estimator actually used, not an assertion that all
its choices follow from one optimal probabilistic model. We derive exact identities,
identify the information that each operation retains or discards, and state the
conditions under which it changes a ranking in the desired direction. Previous
analyses and empirical results are preserved separately; this document adds to them.

\section{Begin with the final score and the two quantities it reconciles}
There are $n$ items and $m$ training bundles. Let $B\in\{0,1\}^{m\times n}$
be their incidence matrix, with $B_{bi}=1$ if item $i$ is in bundle $b$.
Let $S$ be a nonempty observed set, $x\in\{0,1\}^n$ its indicator, and
$c\notin S$ a candidate. Define item frequency $f_i=\sum_b B_{bi}$ and
bundle size $\ell_b=\sum_i B_{bi}$.

We first describe the frozen POG\_dense full recipe. Every $z_h(c,S)$ below is
a standardized channel, defined in Section~\ref{standardize}. Write
\begin{align}
 P_c&=\tfrac12z_{\rm cat}(c,S)+\tfrac14z_{\rm bridge}(c,S)
                         +\tfrac14z_{\rm pop}(c,S),\label{prior}\\
 H_c&=\tfrac12z_{\rm group}(c,S)+\tfrac14z_{\rm csvd}(c,S),\qquad Y_c=P_c+H_c.
 \label{hist}
\end{align}
Here ``group'' is the pair score plus its same-bundle overlap correction, not
the correction alone. The selected constants are $\kappa=0.5$ and
$s=1.4605515925404695$. Define $\rho_c=f_c/(f_c+\kappa)$. Then
\begin{equation}\label{final}
 \boxed{\theta_c=\rho_cY_c+(1-\rho_c)sP_c
       =[s+(1-s)\rho_c]P_c+\rho_c H_c.}
\end{equation}
The subscript $c$ in $P_c,H_c,Y_c,\theta_c$ suppresses their common dependence on $S$.

Why this outer operation? A candidate with no training appearances has no direct
historical support. Its score becomes $sP_c$. As $f_c$ increases, its score
approaches $Y_c$. The fallback is deliberately allowed to share channels with
$Y_c$: the formula adjusts how strongly the extra historical evidence may change
the ranking. $P_c$ is not exclusively textual or independent of history;
category relations and the bridge also use bundles, and popularity is included.

\begin{proposition}[An exact optimization interpretation]\label{objective}
For fixed $Y,P,f$ and $\kappa>0$, the vector in~\eqref{final} is the unique solution of
\begin{equation}\label{energy}
 \min_{t\in\mathbb R^n}\ \frac12\sum_{c=1}^n
       \left[f_c(t_c-Y_c)^2+\kappa(t_c-sP_c)^2\right].
\end{equation}
\end{proposition}
\begin{proof}
The derivative in coordinate $c$ is
$f_c(t_c-Y_c)+\kappa(t_c-sP_c)$, so its zero is
$(f_cY_c+\kappa sP_c)/(f_c+\kappa)$. The Hessian is diagonal with entries
$f_c+\kappa>0$, proving uniqueness, including for cold candidates.
\end{proof}
This says exactly what the convex combination optimizes: agreement with the two
specified scores, not NDCG itself. It does not derive the channels or their
constants. The implementation selects weights, interpolation and shrinkage on
validation rankings using a finite coordinate/grid search. It does not solve a
single jointly convex training objective. When shrinkage is disabled ($\kappa=0$),
the implementation returns $Y$ even for $f_c=0$; the objective with zero coefficient
on that coordinate would not specify a unique answer.

To understand what this compromise achieves, we must now explain the different
evidence in $H$ and $P$, starting with the observations that need no sharing.

\section{Direct support: give a historical bundle a vote, then return its items}
Define $d_i=(f_i+5)^{-1/2}$, $v_b=(\ell_b+1)^{-1/2}$,
$D=\diag(d_i)$ and $V=\diag(v_b)$. The raw pair score is
\begin{equation}\label{pair}
 p=B^\top VBDx,\qquad
 p_c=\sum_{b:c\in b}v_b\sum_{i\in b\cap S}d_i.
\end{equation}
Read the multiplications from right to left. $Dx$ assigns weight to each observed
item. $BDx$ sums those weights within each historical bundle. $V$ discounts the
size of that bundle. $B^\top$ gives its resulting vote to every item it contains.
Thus one matrix multiplication matches the observed set to historical bundles;
the other transfers their support to candidates. No candidate embedding must
compress these individual co-memberships.

For a single recorded co-membership, a rare observed item provides a larger vote
than a common observed item, and a small bundle provides a larger vote than a large
one. These are \emph{per-observation} discounts. They are not a claim that rarity
itself makes a candidate correct or that every bundle has equal total influence.
One matched seed in a size-$\ell$ bundle contributes total candidate mass
$(\ell-1)d_i/\sqrt{\ell+1}$, which still increases with size for $\ell\ge1$.

The asymmetry matters. The frequency factor in~\eqref{pair} is applied to the
\emph{observed item}, not to the candidate. It limits how diagnostic a common clue
is without automatically dividing the candidate's score by its popularity.
Candidate support is treated separately by~\eqref{final}.

An exact alternative view prevents an overstatement about ``debiasing.'' Set
$a_i=\sum_b v_bB_{bi}$. For $a_i>0$, define the size-weighted empirical inclusion
rate $\widehat p_v(c\mid i)=\sum_b v_bB_{bc}B_{bi}/a_i$. Then
\begin{equation}
 p_c=\sum_{i\in S:a_i>0}d_i a_i\,\widehat p_v(c\mid i).
\end{equation}
This combines conditional inclusion rates, but their masses $d_i a_i$ need not be
equal. For equal-sized bundles $a_i=vf_i$, so the mass is
$vf_i/\sqrt{f_i+5}$, approximately proportional to $\sqrt{f_i}$ for frequent
items. If all these empirical conditional rates equal $f_c/m$, the ranking
remains a popularity ranking. The weighting attenuates frequency effects; it
does not remove the independence expectation. The next section's PPMI does a
different normalization for that reason.

\subsection{Why individual matches are not the entire set match}
Let $k_b=(Bx)_b=|b\cap S|$ be the number of observed items in bundle $b$ and
$a_b(S)=(BDx)_b$ its weighted overlap. With $\beta=1.5$, the group score is
\begin{equation}\label{group}
 g_c=\sum_b B_{bc}v_ba_b(S)\{1+\beta[k_b-1]_+\}.
\end{equation}
If only one observed item matches a bundle, the vote is unchanged. If two or more
match, that \emph{same bundle's} vote grows. Support for $c$ from two different
bundles, one matching $a$ and another matching $b$, is distinguished from a bundle
that contains $a,b,c$ together. The correction does not demand that all query
items occur together; it rewards each supported pair of observed items.

Define the triple count $T_{cij}=\sum_bv_bB_{bc}B_{bi}B_{bj}$. For binary queries,
\begin{equation}\label{triple}
 g_c=p_c+\beta\sum_{i<j}x_ix_j(d_i+d_j)T_{cij}.
\end{equation}
To see this, expand $a_b(S)(k_b-1)$; the diagonal terms cancel because $x_i^2=x_i$.
When $k_b=0$, $a_b(S)=0$, so replacing $[k_b-1]_+$ by $k_b-1$ inside this product
does not change it. The coefficients are tied to observed triple counts; they
are not independently fitted tensor entries. This is a quadratic query function,
not an unrestricted model of all subset orders and not a residual against a
pairwise probabilistic null model.

For a fixed query, define
$V_S=\diag(v_b\{1+\beta[k_b-1]_+\})$. Then
\begin{equation}
 g=B^\top V_SB Dx.
\end{equation}
The symmetric middle matrix $B^\top V_SB$ is positive semidefinite: its quadratic
form is a sum of nonnegative weighted squares. The actual asymmetric matrix
$B^\top V_SBD$ is similar to the positive semidefinite matrix
$D^{1/2}B^\top V_SBD^{1/2}$. The query changes the weights of historical bundles
before they are projected back to the item space. These statements include the
diagonal; masking observed candidates occurs later.

There is also an exact normalized-vote interpretation. Let
$w_b(S)=v_ba_b(S)\{1+\beta[k_b-1]_+\}$. If $W(S)=\sum_bw_b(S)>0$,
\begin{equation}
 \frac{g_c}{W(S)}=\sum_b\frac{w_b(S)}{W(S)}B_{bc}.
\end{equation}
This is the probability that $c$ belongs to a historical bundle sampled with those
query-dependent weights. It gives the same raw ranking as $g$. It is not a proof
that these sampling weights are the true posterior probabilities of a future
completion. If no bundle overlaps $S$, all votes are zero.

\begin{proposition}[Item-pair counts cannot recover the group term]\label{loss}
There exist binary training collections with identical $B^\top B$, item
frequencies and bundle sizes, but opposite group rankings for a fixed query.
\end{proposition}
\begin{proof}
Consider the two collections on items $a,b,c,d,e,f$:
\begin{align*}
 \mathcal B_1&=\{\{a,b,c\},\{a,d,e\},\{b,d,f\},\{c,e,f\}\},\\
 \mathcal B_2&=\{\{a,b,d\},\{a,c,e\},\{b,c,f\},\{d,e,f\}\}.
\end{align*}
Every item appears twice. Both contain the same unordered item pairs with the
same multiplicities, hence have identical $B^\top B$. Every bundle has size three.
For $S=\{a,b\}$ the pair scores of $c,d$ are both $1/\sqrt7$. In $\mathcal B_1$,
$c$ receives the factor $1+\beta$ and $d$ does not. In $\mathcal B_2$ the roles
reverse. With the implemented $\beta=1.5$, the scores are
$(0.944911,0.377964)$ and $(0.377964,0.944911)$ respectively.
\end{proof}
Consequently no deterministic function of only $B^\top B$ and identical side
information can recover this distinction in general. The construction was checked
against the actual operators: raw counts, PPMI, the SVD kernel and weighted pair
scores were identical. This is stronger than saying that a low-rank model might
lose information: even the full pair matrix has already lost this information.
It is not a separation from models that also consume bundle membership itself.
Higher-order itemset recommendation has predecessors; the proposition establishes
what this particular term retains, not historical priority.

\section{PPMI and SVD answer different questions about the same pair matrix}
Define $C=\off(B^\top B)$, where $\off$ sets the diagonal to zero. Let
$q_i=\sum_jC_{ij}$ and $T=\sum_{ij}C_{ij}$. For a recorded pair with positive
marginals, define
\begin{equation}\label{ppmi}
 M_{ij}=\pos{\log\frac{C_{ij}T}{q_iq_j}};
 \qquad M_{ij}=0\ \text{for unrecorded pairs}.
\end{equation}
Small numerical stabilizers in the implementation are suppressed here. The raw
co-occurrence and PPMI score vectors are $Cx$ and $Mx$ respectively. The ratio
inside~\eqref{ppmi} compares the empirical probability of an ordered pair with
the product of its pair-endpoint marginals. Thus a frequent observed pair can have
zero PPMI if its count does not exceed that independence baseline. Negative
associations are clipped, not retained as negative evidence.

Here $q_i=\sum_bB_{bi}(\ell_b-1)$ counts \emph{pair endpoints}, whereas $f_i$
counts bundles containing $i$. They differ when bundle sizes vary. PPMI is
therefore not the pair score with a stronger version of the same degree factor:
it uses another sample space, a logarithm and clipping. Both are established
representations of association, with different responses to common and rare pairs.

Now let $V_k\in\mathbb R^{n\times k}$ contain the approximately leading right
singular vectors obtained for $M$. The implementation's transformed rows are
$A=MV_k$, followed by row normalization
$F=RA$, where $R_{ii}=1/(\|A_i\|_2+10^{-8})$. It scores
\begin{equation}\label{svd}
 r_{\rm csvd}=\frac1{|S|}FF^\top x
 =\frac1{|S|}R M V_kV_k^\top M^\top R x.
\end{equation}
This uses the dot product of two projected \emph{association profiles}. It asks
whether candidate and seed relate similarly to other items. It does not directly
look up their reconstructed association $\widehat M_{ci}$.

In an ideal exact SVD $M=U\Sigma V^\top$, the unnormalized kernel is
$U_k\Sigma_k^2U_k^\top$. A rank-$k$ reconstruction instead is
$U_k\Sigma_kV_k^\top$. These are different matrices. At full rank the former
is $MM^\top$, whose entries are $\sum_jM_{cj}M_{ij}$: shared association partners.
For example, two items can have no direct edge yet have identical rows of positive
associations to other items. Their direct PPMI score is zero while their
profile similarity is maximal in the ideal full-rank, normalized case.

There is a further spectral distinction. Since $M$ is symmetric, an exact singular
truncation can be expressed through selected eigenvectors with eigenvalues
$\lambda_j$. Reconstruction retains $\lambda_j$, while the unnormalized profile
kernel uses $\lambda_j^2$. It therefore does not preserve the sign of an
association mode. Selecting the largest singular values also selects by
$|\lambda_j|$, not by a graph notion of smoothness. Calling this channel merely a
smoothed direct-association matrix would hide both changes.

The projection $V_kV_k^\top$ suppresses directions outside its selected subspace.
Exact leading singular vectors minimize
$\|M-MQQ^\top\|_F^2$ over matrices $Q$ with $k$ orthonormal columns. This is a
reconstruction criterion, not a guarantee of relevant-item recovery; the
implementation uses an approximate randomized decomposition. Normalizing rows
also changes the objective's geometry. In particular, shared partners can indicate
substitutes rather than good companions. This explains why keeping direct evidence
alongside the SVD signal is meaningful, and why replacing direct evidence by the
SVD signal is not an equivalent compression.

The full dense recipe uses SVD but selects zero weights for standalone co-occurrence
and PPMI. Their role must not be described as three additional active votes in that
recipe. PPMI remains the input to the selected SVD channel. The transformed-row
interpretation follows the documented SVD convention~\cite{sklearn}; explicit
PPMI-based relational representations have substantial prior history~\cite{levy}.

\section{Attributes share evidence across candidates with little direct support}
The historical operators have an unavoidable boundary: if column $c$ of $B$ is
zero, direct and group support for $c$ vanish, and its PPMI row and SVD vector also
vanish. Descriptors and categories supply another route for ranking that candidate.

Let $G\in\{0,1\}^{n\times a}$ assign each item to one of $a$ nonempty categories,
and let $X\in\mathbb R^{n\times p}$ contain a particular descriptor modality.
Fitting category means solves
\begin{equation}
 \widehat A=\arg\min_A\|X-GA\|_F^2=(G^\top G)^{-1}G^\top X.
\end{equation}
With $Q_G=G(G^\top G)^{-1}G^\top$, the residual is
$X_\perp=(I-Q_G)X$, satisfying $G^\top X_\perp=0$.
The style representation is $E=\normrows(X_\perp)$; its score is
$EE^\top x/|S|$. The subtraction removes the descriptor component explainable
by category means, so this similarity measures deviations from category averages.
The orthogonality statement applies before row normalization; it is not generally
preserved afterwards. It does not imply statistical independence or complete
removal of category information. For multiple modalities the full dense bridge
uses the normalized concatenation of their individually normalized residuals.

Category evidence separately operates at the coarse level. Let $T_G=BG$ count
items of each category within each bundle. The off-diagonal category counts are
$J=\off(T_G^\top T_G)$. Define category occurrence counts
$h_a=\sum_b(T_G)_{ba}$ and $H=\sum_ah_a$. The category relation has entries
$K_{ab}=\pos{\log(J_{ab}H/(h_ah_b))}$ for $a\ne b$, with zero for absent pairs.
For fashion its diagonal is
$K_{aa}=-\tfrac12(1-\delta_a)\mathbf1[\delta_a<0.15]$, where
$\delta_a=(n_{a,\ge2}+0.5)/(n_{a,\ge1}+1)$ and $n_{a,\ge r}$ counts bundles
containing at least $r$ items of category $a$. Its score is
\begin{equation}
 r_{\rm cat}=GKG^\top x.
\end{equation}
Every item in a category shares this category score. This pools observations and
can score an unobserved item pair, but cannot distinguish candidates within the
same category. The diagonal penalizes repeating categories that seldom repeat in
training bundles. These category marginals count occurrences rather than pair
endpoints, so the category log-ratio is not literally the item-PPMI distribution
in~\eqref{ppmi}. It should be called the implemented category association score.

The bridge then connects descriptor residuals to observed partners. Define
\begin{equation}
 L=\off\left(DB^\top\diag((\ell_b+1)^{-1})BD\right),\quad
 N=\normrows(LE),\quad F_{\rm br}=A_{\rm br}(E+\lambda N),
\end{equation}
where $A_{\rm br}$ is the diagonal row-normalization matrix and $\lambda=0.5$.
The graph here discounts both item endpoints and uses inverse bundle size;
it is different from~\eqref{pair}, which uses seed-side frequency and inverse
square-root size. $L$ is symmetric but need not be a Markov matrix or positive
semidefinite after diagonal removal.

The bridge score is $F_{\rm br}F_{\rm br}^\top x/|S|$. Expanding its kernel gives
\begin{equation}\label{bridge}
 F_{\rm br}F_{\rm br}^\top=A_{\rm br}
 (EE^\top+\lambda EN^\top+\lambda NE^\top+\lambda^2NN^\top)A_{\rm br}.
\end{equation}
The middle terms compare an item's descriptor with the aggregated descriptors of
another item's historical partners. Thus two candidates with similar raw
descriptors can differ because their observed relationships differ. For a cold
candidate, its own propagated row is zero, but its descriptor can still match a
warm seed's propagated descriptor. The last term compares their partner summaries.
This is a specified coupling of descriptive similarity and historical relations;
it is not merely a sum of an independent content score and a graph score.

\section{Convert different measurements into comparable ranking margins}\label{standardize}
For channel $h$, let $r_h\in\mathbb R^n$ be its raw scores for $S$. Define
\begin{equation}
 \mu_h=\frac1n\sum_cr_h(c),\quad
 \sigma_h^2=\frac1{n-1}\sum_c(r_h(c)-\mu_h)^2,\quad
 z_h(c)=\frac{r_h(c)-\mu_h}{\sigma_h+10^{-9}}.
\end{equation}
These statistics use the full catalogue, including observed items, before seed
masking; they do not use hidden completion labels. Other queries are not used to
normalize this query. For popularity the raw score is $\log(1+f_c)$.

For $Y_c=\sum_hw_hz_h(c)$ and two candidates $c,d$,
\begin{equation}\label{margin}
 Y_c-Y_d=\sum_h\frac{w_h}{\sigma_h+10^{-9}}[r_h(c)-r_h(d)].
\end{equation}
Hence a channel contributes according to its candidate separation relative to the
spread of its own catalogue scores. This prevents numerical units alone from
determining the result. It does not remove correlation between channels and does
not estimate their uncertainty. In particular, correlated channels can reinforce
the same mistake; a small spread is not proof of high reliability.

Weights are nonnegative, selected on validation rankings. Their sum need not be
one; the dense $Y$ weights sum to $1.75$. The final interpolation in~\eqref{final}
is convex in $Y,sP$, whereas the channel sum is an unnormalized nonnegative sum.
There is another convex interpolation when selecting pair versus group evidence:
$z_{\rm history}=(1-\alpha)z_{\rm pair}+\alpha z_{\rm group}$.
These two channels are standardized separately. It is not generally the same as
standardizing $p+\alpha(g-p)$. Dense selects $\alpha=1$; POG selects $\alpha=0$.

For an equally frequent candidate pair, $\rho_c=\rho_d=\rho$,~\eqref{final} gives
\begin{equation}\label{rankcondition}
 \theta_c-\theta_d=[s+(1-s)\rho](P_c-P_d)+\rho(H_c-H_d).
\end{equation}
If $P$ prefers a correct candidate by $a>0$ and $H$ prefers an incorrect one by
$b>0$, fusion corrects this pair exactly when $[s+(1-s)\rho]a>\rho b$.
When both prefer the correct candidate, both margins reinforce its ranking.
This is a precise mechanism by which complementary evidence improves rankings;
whether correct or incorrect agreements predominate is a question answered by
held-out data. It does not require every channel to be a calibrated probability.

For different frequencies the centering constants also matter. Write
$P_c=\widetilde P_c-A$ and $H_c=\widetilde H_c-J$, where $A,J$ collect the
respective weighted channel means divided by standard deviations and are constant
across candidates within the query. Then
\begin{equation}
 \theta_c=[s+(1-s)\rho_c]\widetilde P_c+\rho_c\widetilde H_c
          -sA-\rho_c[(1-s)A+J].
\end{equation}
The last frequency-dependent term does not cancel in candidate comparisons.
In particular, zero raw history can give a negative standardized history score.
Suppressing that negative contribution for an unseen candidate distinguishes
absence of historical observations from established below-average support, in
the precise sense of this formula. The total adjustment can have either sign.

\section{What the optimization does and does not justify}
Equation~\eqref{energy} provides a complete variational description of the last
step. Frequency determines the cost of departing from $Y$, and $\kappa$ determines
the cost of departing from $sP$. This is a regularization interpretation that
requires no distributional claim. Weighted quadratic penalties are standard;
the value here is identifying their exact role in the implemented score~\cite{boyd}.

A stronger statement, that this is a statistically optimal posterior, would need
additional assumptions about noise and dependence. For illustration, if $Y$ and
$sP$ were unbiased estimators of the same scalar with error variances $v_Y,v_P$
and error covariance $u$, their minimum-MSE mixture weight on $Y$ would be
\begin{equation}
 \rho^*=\frac{v_P-u}{v_Y+v_P-2u},
\end{equation}
provided the denominator is positive (and clipped to $[0,1]$ for a convex mixture).
The rational weight $f/(f+\kappa)$ follows in the special case $u=0$,
$v_Y=a/f$, $v_P=b$ and $\kappa=a/b$. The actual scores share channels, so this
special case is not established. The dense frozen recipe's fitted raw variance
slope is negative, about $-2.5281$; its selected $\kappa=0.5$ is a validation
choice, not evidence that a positive $1/f$ variance law was recovered. We therefore
base the argument on the exact optimization and measured rankings, not on a
claimed empirically verified Bayesian data-generating model.

\section{The resulting scientific argument and its evidential boundary}
The method can be described precisely as follows:
\begin{quote}
Estimate completion support without forcing different relations into one
similarity matrix. Preserve direct weighted co-membership and same-bundle joint
support; share association patterns through a projected PPMI profile, category
relations and graph-conditioned descriptors. Combine their catalogue-relative
ranking margins, then regularize the combined score toward shared evidence more
strongly when a candidate has fewer observations.
\end{quote}
Three mathematical results support this account. First, identical pair counts can
conceal different same-bundle evidence (Proposition~\ref{loss}). Second, the SVD
and bridge kernels explicitly introduce shared-partner and descriptor--partner
comparisons (Equations~\eqref{svd} and~\eqref{bridge}); these are distinct relations,
not repetitions of direct counts. Third, final fusion solves a strictly convex
support-weighted reconciliation (Proposition~\ref{objective}).

The recipe does not rely on every possible signal. On POG, the selected history
term is only $H_c=0.25z_{\rm pair}$, and
$P_c=0.5z_{\rm cat}+0.25z_{\rm style}+0.25z_{\rm bridge}$, with
$\kappa=0.5$ and $s\simeq1.0424024$. Both incidence and SVD are unselected.
The common argument is retaining direct evidence while sharing information and
adjusting its influence to support. Joint support specifically strengthens the
dense recipe; it cannot explain the POG result where its coefficient is zero.

The saved mechanism studies report Recall@20 for history, shared, combined and
full as $0.375379,0.112917,0.390928,0.395223$ on POG\_dense and
$0.030875,0.044542,0.051458,0.054458$ on POG. The standalone families are fitted
separately and the shared family includes graph information; these are empirical
complementarity observations, not independent measurements of each operator's
causal effect. The dense full versus matched-frequency comparison improves
NDCG@20 from $0.279345$ to $0.287876$ with the group term enabled; its fallback
scales differ slightly, so this is not a perfectly isolated intervention.
The recall contrast does not pass the study's multiple-testing correction.

The reported superiority to published neural and evaluated fixed-score baselines
is an empirical claim to state with its actual inputs, splits and metrics. The
algebra above explains structural capabilities and conditional ranking effects;
it is not a universal performance theorem or a claim that previous neural models
cannot represent these relations. Related neighborhood, itemset, representation
and confidence-adaptive methods remain relevant; the preserved prior-art review
records those connections.

\appendix
\section{Evidence and reproduction}
The primary scope is binary, unit-weight bundles and the full-ranking research
recipes. Weighted retail profiles need their own derivation. The checks use the
actual graph and fusion operators on explicit synthetic inputs and record their
source hashes. They do not fit or tune against test targets.

\begin{itemize}
\item New derivation checks: \path{scripts/check_scoring_principle.py}; output
\path{results/algebraic_diagnostics/scoring_principle.json}.
They check identical-pair corpora, SVD projection and Gram identities,
query-dependent kernels, normalized historical voting, category least squares,
bridge expansion and stationarity of the exact final objective.
\item Earlier identities and counterexample:
\path{docs/algebraic_mechanism_analysis.md} and
\path{scripts/check_higher_order_structure.py}.
\item Empirical complementarity:
\path{results/algebraic_diagnostics/evidence_overlap.json}, generated by
\path{scripts/analyze_evidence_overlap.py} from saved query outcomes.
\item Recipe and metric evidence: each dataset's \path{frozen_recipes.json},
\path{summary.json} and \path{queries.json} under
\path{results/publication_mechanisms/}.
\item Full source correspondence: pair, group, PPMI, SVD, descriptor centering and
bridge in \path{bundlefuse/operators.py}; query feature aggregation and
normalization in \path{bundlefuse/selected_signals.py}; weighted sums,
selection and backoff in \path{bundlefuse/fusion.py}; modality assembly in
\path{bundlefuse/model.py}.
\item Earlier publication arguments and literature review are retained;
\path{docs/analysis_evidence_index.md} explains their relationship to this note.
\end{itemize}

\begin{thebibliography}{9}
\bibitem{sklearn} Scikit-learn. Truncated singular value decomposition and latent
semantic analysis, User Guide, Section 2.5.3.
\url{https://scikit-learn.org/stable/modules/decomposition.html#truncated-singular-value-decomposition-and-latent-semantic-analysis}.
\bibitem{levy} O. Levy and Y. Goldberg. Linguistic Regularities in Sparse and
Explicit Word Representations. CoNLL, 2014.
\url{https://aclanthology.org/W14-1618/}.
\bibitem{boyd} S. Boyd and L. Vandenberghe. Convex Optimization. Cambridge
University Press, 2004. Chapter 6.
\url{https://web.stanford.edu/~boyd/cvxbook/}.
\end{thebibliography}
\end{document}
